\section{Comparison and support tools under condition (C)}
\label{sec:support-tools}

Corollary~\ref{ext:cor:condition-C} has supplied condition (C) from
simple extended-functional comparison. We first complete the other
route, from the canonical Cuntz-semigroup map. The remaining lemmas turn
(C) into the hereditary supports, matrix cones, and bounded witnesses
used by the common absorption construction. They apply without
simplicity; exactness is needed for the trace comparison, and
separability enters the central-sequence constructions.

\subsection{The consequences of the canonical Cu map}

\begin{lemma}[Cuntz regularity]
\label{lem:C}
Under the hypothesis of Corollary~\ref{cor:canonical}, the semigroup
$\Cu(A)$ satisfies condition \textup{(C)}, namely
Equations~\eqref{eq:pre:au}--\eqref{eq:pre:div}.
\end{lemma}

\begin{proof}
The algebra \(A\otimes\Z\) and its stabilization absorb \(\Z\).
R{\o}rdam's Theorem~4.5~\cite{Rordam2004} makes their \(W\)-semigroup
almost unperforated. Transport through \(\Cu(\iota)\) proves
\eqref{eq:pre:au}.

For an integer $N\ge1$, R\o rdam's divisibility lemma
\cite[Lemma~4.2]{Rordam2004} supplies $z\in\Z_+$ such that
\[
 N[z]\le[1_{\Z}]\le(N+1)[z].
\]

For \(x=[b]\), tensor these comparisons with \(b\). Comparison is
preserved by tensoring with a fixed positive element
\cite[Lemma~4.1]{Rordam2004}: tensor the witnesses and use a bounded
approximate unit on the other factor.
Thus
\[
 N[b\otimes z]\le[b\otimes1_{\Z}]
       \le(N+1)[b\otimes z].
\]
The middle class is exactly \(\Cu(\iota)(x)\). Pulling back the
outside class proves \eqref{eq:pre:div}.
\end{proof}

\subsection{Comparison from a rank gap}

The later construction will compare a defect with a matrix-cone weight
by proving a uniform gap between their ranks. The following form allows
traces that vanish on an ideal or take infinite values. Exactness turns
the rank inequalities into functional inequalities, and almost
unperforation supplies comparison through
Lemma~\ref{lem:comparison-convention}.

\begin{lemma}[Comparison from a uniform rank gap]
\label{lem:G}
Suppose \(A\) is exact and satisfies \eqref{eq:pre:au}. If \(v,w\in A_+\)
satisfy
\[
 d_\tau(v)\le\tfrac12d_\tau(w)
\]
for every lower semicontinuous extended trace on \(A\), then
\(v\precsim w\).
\end{lemma}

\begin{proof}
Let $\varphi\in\F(\Cu(A))$. By
Equation~\eqref{ext:eq:trace-representation}, it is represented by an
extended trace on $A\otimes\K$. Restricting that trace to the original
corner and applying the hypothesis gives
\[
 \varphi([v])\le\tfrac12\varphi([w])\le\varphi([w]).
\]
The inequality is strict whenever $0<\varphi([w])<\infty$.
Since $\Cu(A)$ is almost unperforated,
Lemma~\ref{lem:comparison-convention} applies and gives $[v]\le[w]$.
\end{proof}

\subsection{Placing divisible matrix cones}

Almost divisibility provides equal Cuntz classes below a prescribed
positive element. Cutting these classes lets us place the corresponding
matrix cone inside its hereditary algebra. The two forms below retain,
respectively, norm visibility in representations and a prescribed
fraction of the rank for bounded traces.

\begin{lemma}[Cone placement]
\label{lem:placement}
Assume \eqref{eq:pre:au}--\eqref{eq:pre:div}.
\begin{enumerate}
\item Let \(B\) be a hereditary subalgebra of a matrix algebra over
\(A\otimes\K\), let \(b\in B_+\) be a contraction, and let \(N\ge1\)
and \(\epsilon>0\). There is a c.p.c.\ order-zero map
\(\psi:M_N\to\her_B(b)\) such that
\[
 \|\sigma(\psi(1))\|=1
 \quad\text{whenever}\quad
 \sigma((b-\epsilon)_+)\ne0.
\]
Its \(N\) diagonal weights are orthogonal and equivalent in
\(\her_B(b)\). Each diagonal weight, cut at \(1/2\), dominates a
positive class \(u_0\) for which
\[
 [(b-\epsilon)_+]\le(N+1)u_0.
\]
\item Let \(C\) be isomorphic to a finite direct sum of hereditary
subalgebras as above. For \(b\in C_+\) a contraction, \(k\ge1\), and
\(0<q\le1\), there is a c.p.c.\ order-zero map
\(\psi:M_k\to\her_C(b)\), with weight \(t=\psi(1)\), such that
\begin{equation}
\label{eq:pre:P}
 d_\sigma(t)\le q\,d_\sigma(b),\qquad
 d_\sigma((t-\tfrac12)_+)
       \ge(q/3)d_\sigma((b-\tfrac12)_+)
\end{equation}
for every bounded positive trace \(\sigma\) on \(C\).
\end{enumerate}
\end{lemma}

\begin{proof}
Choose \(u=[r]\) with \(N[r]\le[b]\le(N+1)[r]\).
By \eqref{eq:pre:cuts}, a sufficiently small \(\eta>0\) gives
\[
 [(b-\epsilon)_+]\le(N+1)[(r-2\eta)_+].
\]
Lemma~\ref{lem:pre:cut-placement} places the hereditary algebra of
\(\operatorname{diag}((r-\eta)_+,\ldots,(r-\eta)_+)\), with \(N\)
diagonal entries, inside \(\her_B(b)\). In this hereditary matrix
algebra take the cone \(e_{ij}\mapsto e_{ij}\otimes g(r)\), where
\(0\le g\le1\), \(g=0\) on \([0,\eta]\), and \(g=1\) on
\([2\eta,\infty)\). Transport it by the polar isometry.
Every diagonal acts as one on the corresponding transported
\((r-2\eta)_+\), proving both its cut comparison and the norm-one
assertion. The off-diagonal entries give equivalence of the
orthogonal weights.

For the second assertion do this in each summand, with
\(N=\lceil k/q\rceil\) and \(\epsilon=1/2\), and restrict to the first
\(k\) slots. The complete \(N\)-slot weight lies in \(\her(b)\),
so equivalence and additivity of rank give
\[
 d_\sigma(t)\le(k/N)d_\sigma(b)\le q\,d_\sigma(b).
\]
The retained narrow classes and the cut comparison give
\[
 d_\sigma((t-\tfrac12)_+)
 \ge\frac{k}{N+1}d_\sigma((b-\tfrac12)_+)
 \ge(q/3)d_\sigma((b-\tfrac12)_+).
\]
All comparisons take place in the indicated hereditary algebra,
by Lemma~\ref{lem:pre:cut-placement}. Thus arbitrary bounded traces
on that algebra may be used, whether or not they extend boundedly
to \(A\). Summing the finitely many summands proves the claim.
\end{proof}

\subsection{A support compression}

To make an error operator small on a hereditary support, we compress its
complement below a chosen threshold and take the positive part. The
following observation gives the resulting bound on the support, without
requiring the error operator to commute with the compressing element.

\begin{lemma}
\label{lem:support}
Let \(K=K^*\) be a bounded operator on a Hilbert space, let \(s\ge0\)
be bounded, and let \(\lambda\in\R\). On the support of
\[
 b=s[s(\lambda-K)s]_+s
\]
the compression of \(K\) is at most \(\lambda\).
If \(b=0\), then \(s(\lambda-K)s\le0\).
\end{lemma}

\begin{proof}
Put \(H=s(\lambda-K)s\). For a vector \(\xi\) in the positive spectral
space of \(H\),
\[
 \langle(\lambda-K)s\xi,s\xi\rangle
       =\langle H\xi,\xi\rangle\ge0.
\]
The range closure of \(sH_+^{1/2}\) is the support space of \(b\).
It is contained in the closure of the images \(s\xi\) just used.
Continuity of the quadratic form proves the compression bound.

If \(b=0\), then \(H_+^{1/2}s=0\). Since \(H_+\) is supported on
\(p_s\), and the range of \(s\) is dense in its support space, this implies
\(H_+=0\), as required.
\end{proof}

\subsection{Ultrapowers and bounded diagonalization}

Fix a free ultrafilter \(\omega\) on \(\N\). We write
\[
 D_\omega=\ell^\infty(\N,D)/
   \{(x_n):\lim_{n\to\omega}\|x_n\|=0\},
\]
and identify \(D\) with its constant copy. For separable \(A\),
\[
 \Ann(A)=\{x\in A_\omega\cap A':xa=ax=0
             \text{ for every }a\in A\}.
\]
An approximate unit of \(A\) gives the unit in
\((A_\omega\cap A')/\Ann(A)\).
We use c.p.c.\ lifting for separable nuclear domains
\cite[Theorem~3.10 and Corollary~3.11]{CE1976}; in particular maps from a matrix algebra can be lifted
to c.p.c.\ coordinate maps.

\begin{lemma}[Bounded diagonalization]
\label{lem:diagonal}
Suppose a countable family of norm-continuous algebraic and
functional-calculus conditions, with fixed coefficients in a norm
ultrapower, is approximately satisfiable on every finite subfamily.
Suppose each unknown has a fixed individual norm bound in these
finite solutions. Then the conditions can be satisfied exactly in
the ultrapower. Additional countably many commutator conditions and
conditions after multiplication by a fixed countable subset of the
algebra may be imposed in the same way.
\end{lemma}

\begin{proof}
For each \(j\), choose an ultrapower solution of the first \(j\)
conditions to accuracy \(1/j\), and choose bounded coordinate
representatives. Slightly enlarging each fixed bound once, and
truncating representatives, gives bounds independent of \(j\).
There is an \(\omega\)-large set on which these first \(j\)
coordinate conditions hold to accuracy \(2/j\). Replace these sets
by decreasing finite intersections, also intersecting the \(j\)-th
set with \(\{n\ge j\}\). They have empty intersection. At coordinate
\(n\), use the solution of the largest applicable level \(j\le n\),
with an arbitrary bounded choice when no level applies.
The chosen level tends to infinity along \(\omega\); each fixed
condition therefore has error tending to zero. The bounds ensure
that every unknown represents an element of the norm ultrapower.
Positivity or complete positivity can be preserved by choosing
positive or completely positive lifts at each level.
\end{proof}

We will also use the following reindexing observation. If a central
sequence is to commute with finitely many already chosen coordinate
values, then, at each output coordinate, those values are ordinary
elements of \(A\). Centrality supplies the required good set of old
indices. It can be intersected with finitely many other good sets
before choosing the new index. Any uniform coordinate predicate
retained in this selection continues to hold. In the nonunital
quotient, algebraic and unit defects are tested after multiplication
by a dense subset of \(A\); boundedness then extends the resulting
relations to all of \(A\).
